\begingroup
\renewcommand{\arraystretch}{1.45}
\Needspace{46\baselineskip}
\InterpolationArticle{37}

In den Artikeln 30---36 haben wir angenommen, dass $\sin\mu a$ nicht $=0$ sei: Es bleibt daher, dieselben Untersuchungen für den Fall wieder aufzunehmen, in dem $\sin\mu a=0$ ist. Hier werden wir sogleich $a=0$ oder $a=\frac{180^{\mathord{0}}}{\mu}$ voraussetzen.
\SourceLeaf{329}
\par Zunächst sei $X$ eine Funktion der Form
\[
\begin{aligned}
&\gamma+\gamma'\cos x+\gamma''\cos2x+\mathrm{usw.}+\gamma^n\cos nx
\end{aligned}
\]
sodass die Anzahl der unbekannten Koeffizienten $=1+n$ ist. Wir wissen bereits, dass von allen vorgelegten Werten der Funktion $X$ diejenigen verworfen werden müssen, die einem Argumentwert größer als $180^{\mathord{0}}$ entsprechen: Hier sind daher vier Fälle zu unterscheiden:

1) Wenn $\mu$ gerade und $a=0$ ist, so wird $180^{\mathord{0}}$ der $\frac12\mu+1^{\mathrm{te}}$ Wert von $x$ sein; da die folgenden also verworfen werden müssen, bleiben $\frac12\mu+1$ Werte übrig. Daher muss $n=\frac12\mu$ sein.

2) Es sei $\mu$ gerade und $a=\frac{180^{\mathord{0}}}{\mu}$. Der $\frac12\mu^{\mathrm{te}}$ Wert von $x$ wird $180^{\mathord{0}}-\frac{180^{\mathord{0}}}{\mu}$ sein; die folgenden, die größer als $180^{\mathord{0}}$ werden, sind zu verwerfen. Daher muss $n=\frac12\mu-1$ sein.

3) Wenn $\mu$ ungerade und $a=0$ ist, so ergibt sich für den Wert $\frac12\mu+\frac12^{\mathrm{te}}=180^{\mathord{0}}-\frac{180^{\mathord{0}}}{\mu}$, und

4) Wenn $\mu$ ungerade und $a=\frac{180^{\mathord{0}}}{\mu}$ ist, so ergibt sich für den Wert $\frac12\mu+\frac12^{\mathrm{te}}=180^{\mathord{0}}$: Die folgenden müssen in beiden Fällen verworfen werden, sodass $n=\frac12\mu-\frac12$ sein wird.

\Needspace{32\baselineskip}
Da nun die in den vorigen Artikeln angewendete Methode auf den gegenwärtigen Fall, in dem zuvor ein Teil der gegebenen Werte von der vollständigen Periode abgeschnitten werden müsste, nicht ohne gewisse Umwege angewendet werden könnte, ziehen wir die folgende Methode vor.

Wird nach den Vorschriften der Artikel 20, 22 eine Funktion der Form
\[
\begin{aligned}
&\gamma+\gamma'\cos x+\gamma''\cos2x+\mathrm{usw.}+\gamma^m\cos mx\\[5pt]
&{}+\delta'\sin x+\delta''\sin2x+\mathrm{usw.}+\delta^m\sin mx
\end{aligned}
\]
gesucht, die allen $\mu$ gegebenen Werten genügt und in der $m=\frac12\mu-\frac12$ oder $=\frac12\mu$ ist, je nachdem, ob $\mu$ ungerade oder gerade ist, so werden die Koeffizienten $\delta'$, $\delta''$, $\delta'''$ usw. von selbst $=0$. Denn bei einem Ausdruck wie
\[
\begin{aligned}
&A\sin\lambda a+B\sin\lambda b+C\sin\lambda c+D\sin\lambda d+\mathrm{usw.}
\end{aligned}
\]
ist leicht ersichtlich, dass entweder der erste Teil $=0$ wird, der letzte $=-B\sin\lambda b$, der vorletzte $=-C\sin\lambda c$, der drittletzte $=-D\sin\lambda d$ usw., nämlich wenn $a=0$ ist; oder der letzte $=-A\sin\lambda a$, der vorletzte $=-B\sin\lambda b$, der drittletzte $=-C\sin\lambda c$ usw., wenn $a=\frac{180^{\mathord{0}}}{\mu}$ ist, da für solche Werte von $x$, von denen der eine den anderen zu $360^{\mathord{0}}$ ergänzt, die Werte der Funktion $X$ gleich sind. \Needspace{46\baselineskip}
Daher wird die Funktion
\[
\begin{aligned}
&X'=\gamma+\gamma'\cos x+\gamma''\cos2x+\mathrm{usw,}+\gamma^m\cos mx
\end{aligned}
\]
\SourceLeaf{330}
in der die Koeffizienten $\gamma$, $\gamma'$, $\gamma''$ usw. durch die Formeln
\[
\begin{aligned}
&\gamma = \frac{1}{\mu}(A+B+C+D+\mathrm{usw.})\\[5pt]
&\gamma' = \frac{2}{\mu}(A\cos a+B\cos b+C\cos c+D\cos d+\mathrm{usw.})\\[5pt]
&\gamma'' = \frac{2}{\mu}(A\cos 2a+B\cos 2b+C\cos 2c+D\cos 2d+\mathrm{usw.})
\end{aligned}
\]
usw., der letzte aber, wenn $\mu$ gerade und somit $m=\frac12\mu$ ist, durch diese
\[
\begin{aligned}
&\gamma^m=\frac{1}{\mu}(A\cos ma+B\cos mb+C\cos mc+D\cos md+\mathrm{usw.})
\end{aligned}
\]
bestimmt werden, notwendig mit der Funktion $X$ identisch sein, sofern diese nicht von höherem Grad ist, als wir oben festgelegt haben. Denn in den Fällen $3$ und $4$ ist $X$ von der $\frac12\mu-\frac12{}^{\mathrm{ten}}$ Ordnung, d. h. von derselben wie $X'$ und daher nach Artikel 24 mit $X'$ identisch. Im ersten Fall ist $X$ von derselben Ordnung wie $X'$ und im zweiten von keiner höheren; daher werden in beiden Fällen alle Terme wenigstens bis zur $\frac12\mu-1^{\mathrm{ten}}$ Ordnung in beiden Funktionen übereinstimmen (Artikel 24). Dass auch die Terme von der $\frac12\mu^{\mathrm{ten}}$ Ordnung in diesen Funktionen übereinstimmen müssen, geht nach demselben Artikel 24 daraus hervor, dass in $X$ die Bedingungsgleichung $K\sin ma=L\cos ma$ gilt; es wird nämlich $L=0$ und im ersten Fall $\sin ma=0$, im zweiten aber, wo $X$ nur bis zur Ordnung $\frac12\mu-1$ ansteigt, $K=0$. Im zweiten Fall scheint übrigens $X'$ von höherer Ordnung zu sein als $X$; aber in diesem Fall verschwindet auch der Term von der $\frac12\mu^{\mathrm{ten}}$ Ordnung in $X'$, da
\[
\begin{aligned}
&\gamma^{\frac12\mu}=\frac{1}{\mu}(A\cos 90^{{\mathord{0}}}+B\cos 270^{{\mathord{0}}}+C\cos 450^{{\mathord{0}}}+D\cos 630^{{\mathord{0}}}+\mathrm{usw.})=0
\end{aligned}
\]
sodass auch in diesem Fall $X'$ tatsächlich von der $m-1^{\mathrm{ten}}$ Ordnung oder der $\frac12\mu-1^{\mathrm{ten}}$ ist.

\Needspace{46\baselineskip}
\InterpolationArticle{38}

Wenn die Funktion $X$ nicht mit dem Term $\cos nx$ abbricht, sondern weiterläuft, so wird, wenn die folgenden Terme durch \hfill\break{}
$\alpha^{n+1}\cos(n+1)x+\alpha^{n+2}\cos(n+2)x+\mathrm{usw.}$ bezeichnet werden, nach den Artikeln 21, 23
\[
\begin{aligned}
&\gamma = \alpha\pm\alpha^{\mu}+\alpha^{2\mu}\pm\mathrm{usw.}\\[5pt]
&\gamma' = \alpha'\pm\alpha^{\mu-1}\pm\alpha^{\mu+1}+\alpha^{2\mu-1}+\alpha^{2\mu+1}\pm\mathrm{usw.}\\[5pt]
&\gamma'' = \alpha''\pm\alpha^{\mu-2}\pm\alpha^{\mu+2}+\alpha^{2\mu-2}+\alpha^{2\mu+2}\pm\mathrm{usw.}\\[5pt]
&\gamma''' = \alpha'''\pm\alpha^{\mu-3}\pm\alpha^{\mu+3}+\alpha^{2\mu-3}+\alpha^{2\mu+3}\pm\mathrm{usw.}
\end{aligned}
\]
und so fort bis zum letzten Koeffizienten $\gamma^m$, wenn $\mu$ ungerade ist, oder bis zum vorletzten $\gamma^{m-1}$, wenn $\mu$ gerade ist; das untere Zeichen gilt hier, sooft $a=\frac{180^{{\mathord{0}}}}{\mu}$ ist, also
\SourceLeaf{331}
in den Fällen $2$ und $4$, das obere in den Fällen $1$ und $2$; schließlich wird für den letzten Koeffizienten im ersten Fall
\[
\begin{aligned}
&\gamma^m=\alpha^m+\alpha^{3m}+\alpha^{5m}+\mathrm{usw.}
\end{aligned}
\]
im zweiten Fall dagegen $\gamma^m=0$ sein. Diese Gleichungen zeigen, inwieweit es zulässig sein könnte, den Unterschied zwischen den Funktionen $X$ und $X'$ zu vernachlässigen. Diese letztere Funktion wird unter allen, die den $\mu$ vorgelegten Werten genügen, die einfachste sein; alle diese werden in der Form $X'+Y\sin(\frac12\mu x-\frac12\mu a)$ enthalten sein, die in den Fällen $1$ und $3$ auf $X'+Y\sin\frac12\mu x$, in den Fällen $2$ und $4$ dagegen auf $X'+Y\cos\frac12\mu x$ zurückgeht. Soll dieser Ausdruck aber dieselbe Form wie $X$ haben, d. h. frei von Sinusgliedern sein, so muss offenbar $Y$ sein
\[
\begin{aligned}
&\text{im ersten Fall von der Form} \quad  g\sin x+g'\sin 2x+g''\sin 3x+\mathrm{usw.}\\[5pt]
&\text{im zweiten Fall von der Form} \quad  g+g'\cos x+g''\cos 2x+\mathrm{usw.}\\[5pt]
&\text{im dritten Fall von der Form} \quad  g\sin\frac12 x+g'\sin\frac32 x+g''\sin\frac52 x+\mathrm{usw.}\\[5pt]
&\text{im vierten Fall von der Form} \quad  g\cos\frac12 x+g'\cos\frac32 x+g''\cos\frac52 x+\mathrm{usw.}
\end{aligned}
\]
\Needspace{36\baselineskip}
Und allgemeiner: Bezeichnet $X''$ irgendeine der Funktion $X$ ähnliche Funktion, die den $\mu$ vorgelegten Werten genügt, so werden alle derartigen Formen in der Formel $X''+Y\sin\frac12\mu x$ oder $X''+Y\cos\frac12\mu x$ enthalten sein, wobei $Y$ eine unbestimmte Funktion von der eben angegebenen Form bezeichnet. Dies lässt sich so ausführen, dass die Funktion auf die Ordnung $\frac12\mu$, $\frac12\mu-1$, $\frac12\mu-\frac12$, $\frac12\mu-\frac12$ herabgesetzt hervorgeht und offenbar mit $X'$ identisch sein wird. Die allgemeine Regel zur Herabsetzung einer solchen Funktion $X''$ auf $X'$ ergibt sich leicht aus Artikel 24. Für jeden Term $L\cos\lambda x$ in $X''$ muss man in $X'$, indem man $\lambda=k\mu\pm\lambda'$ setzt, sodass $\lambda'$ nicht größer als $\frac12\mu$ ist, den Term $\pm L\cos\lambda'x$ einsetzen; dabei ist das untere Zeichen zu nehmen, sooft zugleich $a=\frac{180^{{\mathord{0}}}}{\mu}$ und $k$ gerade ist, das obere in den übrigen Fällen; schließlich muss, sooft im zweiten Fall, d. h. für $a=\frac{180^{{\mathord{0}}}}{\mu}$ und einen geraden Wert von $\mu$, $\lambda'=\frac12\mu$ wird, für $L\cos\lambda x$ sogleich $0$ in $X'$ gesetzt, also jener Term ganz weggelassen werden.

\endgroup
